\section {Five pirates (Logic)}

Five pirates find a sack with 100 gold coins and want to divide them.
The pirates are ranked by seniority: the most senior is \#5, then \#4 and so on down to the most junior, \#1.
The most senior pirate proposes how to divide the coins and then all the pirates, including him, vote on the proposal.
If at least half of the votes accept it, the coins are divided accordingly and they are done.
Otherwise the pirate who made the proposal is shot, and the next in line makes a proposal, and so on.

Every pirate is perfectly rational. His first priority is to stay alive, his second is to get as many coins as possible,
and his third is blood: if a vote makes no difference to him otherwise, he votes against, hoping for a kill.

What should pirate \#5 propose to stay alive and keep as many coins as possible?

\subsection{Solution}

Pirate \#5 gives one coin to pirate \#3 and one coin to pirate \#1 and keeps the other 98.

To see why, solve the problem backwards, starting with fewer pirates.

\textbf{Two pirates.} \#2 proposes to keep all 100 coins. He votes for it, which is already half of the votes, so it passes.
\#1 gets nothing.

\textbf{Three pirates.} \#3 needs two votes. He knows that if he is shot, \#1 gets nothing.
So he offers \#1 a single coin, which is better for \#1 than what he would get otherwise, and \#1 votes for it.
Proposal: \#3 gets 99, \#2 gets 0, \#1 gets 1.

\textbf{Four pirates.} \#4 needs two votes out of four, his own and one more.
If he is shot, \#2 gets nothing, so a single coin buys his vote.
Proposal: \#4 gets 99, \#3 gets 0, \#2 gets 1, \#1 gets 0.

\textbf{Five pirates.} \#5 needs three votes out of five, his own and two more.
If he is shot, \#3 and \#1 get nothing, so a single coin buys each of their votes.
Proposal: \#5 gets 98, \#4 gets 0, \#3 gets 1, \#2 gets 0, \#1 gets 1.

Note that giving a pirate exactly what he would get anyway is not enough to buy his vote, since he would then vote against for the blood.
